\begin{ouverture}[geometrie]{13}{Le théorème de Pythagore}{Partie 4 · Géométrie et espace}
{Triangle rectangle · Longueurs · Réciproque}
\objectif{Calculer un côté d’un triangle rectangle.}
\objectif{Démontrer qu’un triangle est rectangle ou ne l’est pas.}
\objectif{Rédiger et contrôler une solution géométrique.}
\prerequis{Carrés, racines carrées et triangles.}
\end{ouverture}
\pagestyle{otMain}\essentiel\fontsize{11}{15}\selectfont
\phantomsection\label{cours-13-1}
\essSection{1}{Calculer dans un triangle rectangle}
\begin{propriete}[Théorème de Pythagore]Si ABC est rectangle en A, alors $BC^2=AB^2+AC^2$. Le côté opposé à l’angle droit, BC, est l’hypoténuse : c’est le plus long côté. Le théorème relie les aires des carrés construits sur les trois côtés.\end{propriete}
\label{fig-cours-13-1-cours-pythagore}\figureGeo{cours-pythagore}

\begin{exemple}Si $AB=3$ et $AC=4$ cm, alors $BC^2=9+16=25$, donc $BC=\sqrt{25}=5$ cm. Si $BC=13$ et $AB=5$ cm, alors $AC^2=13^2-5^2=144$, donc $AC=12$ cm.\end{exemple}
\begin{attention}On additionne les carrés pour chercher l’hypoténuse ; on soustrait pour chercher un côté de l’angle droit. Écrire d’abord l’égalité avec les noms des points aide à éviter cette confusion.\end{attention}
\clearpage
\phantomsection\label{cours-13-2}
\essSection{2}{Reconnaître et prouver}
\begin{propriete}[Réciproque]Si, dans un triangle, le carré du plus long côté est égal à la somme des carrés des deux autres, alors ce triangle est rectangle. L’angle droit est opposé au plus long côté.\end{propriete}
\begin{savoirfaire}[Prouver ou réfuter]Identifier le plus long côté. Calculer séparément son carré et la somme des deux autres carrés. Comparer : une égalité permet d’utiliser la réciproque ; une inégalité montre que le triangle n’est pas rectangle, car un triangle rectangle respecterait Pythagore.\end{savoirfaire}
\begin{exemple}Pour des côtés 6, 8, 10, $10^2=100=36+64$ : le triangle est rectangle. Pour 6, 8, 11, $11^2=121\ne100$ : il ne l’est pas. Un dessin qui paraît rectangle ne suffit pas.\end{exemple}
\begin{savoirfaire}[Rédiger]Donner le triangle et l’hypothèse d’angle droit, citer le théorème, écrire l’égalité, calculer, prendre la racine positive et conclure avec l’unité. Pour une réciproque, annoncer la comparaison puis le sommet de l’angle droit. Garder la valeur exacte avant l’arrondi final.\end{savoirfaire}

\clearpage\rubrique{Découverte}
\begin{blocLivre}{ex-13-1}\begin{exercice}[1]{Chercher l’hypoténuse}
ABC est rectangle en A, $AB=6$ cm et $AC=8$ cm. Calculer BC.
\label{fig-ex-13-1}\figureGeo{pyth-01}
\end{exercice}\end{blocLivre}
\begin{blocLivre}{ex-13-2}\begin{exercice}[1]{Chercher un petit côté}
ABC est rectangle en A, $AB=5$ cm et $BC=13$ cm. Calculer AC.
\label{fig-ex-13-2}\figureGeo{pyth-02}
\end{exercice}\end{blocLivre}
\begin{blocLivre}{ex-13-3}\begin{exercice}[1]{Une racine non entière}
ABC est rectangle en A avec $AB=7$ cm et $AC=9$ cm. Donner BC exactement puis au dixième.
\label{fig-ex-13-3}\figureGeo{pyth-03}
\end{exercice}\end{blocLivre}
\begin{blocLivre}{ex-13-4}\begin{exercice}[1]{L’hypoténuse donnée}
ABC est rectangle en A, $BC=17$ cm et $AB=15$ cm. Calculer AC.
\label{fig-ex-13-4}\figureGeo{pyth-04}
\end{exercice}\end{blocLivre}
\begin{blocLivre}{ex-13-5}\begin{exercice}[1]{Une réciproque}
Un triangle ABC a pour côtés $AB=9$, $AC=12$, $BC=15$ cm. Est-il rectangle ?
\label{fig-ex-13-5}\figureGeo{pyth-05}
\end{exercice}\end{blocLivre}

\rubrique{Entraînement}
\begin{blocLivre}{ex-13-6}\begin{exercice}[2]{Un triangle non rectangle}
Un triangle ABC a pour côtés $AB=4$, $AC=5$, $BC=7$ cm. Est-il rectangle ?
\label{fig-ex-13-6}\figureGeo{pyth-06}
\end{exercice}\end{blocLivre}
\begin{blocLivre}{ex-13-7}\begin{exercice}[2]{Diagonale d’un rectangle}
Un rectangle mesure 8 cm sur 6 cm. Calculer la longueur de sa diagonale.
\label{fig-ex-13-7}\figureGeo{pyth-07}
\end{exercice}\end{blocLivre}
\begin{blocLivre}{ex-13-8}\begin{exercice}[2]{Une échelle dans un modèle}
Une échelle de 2,5 m repose contre un mur perpendiculaire au sol. Son pied est à 0,8 m du mur. À quelle hauteur arrive-t-elle, au centimètre près ?
\label{fig-ex-13-8}\figureGeo{pyth-08}
\end{exercice}\end{blocLivre}
\begin{blocLivre}{ex-13-9}\begin{exercice}[2]{Une distance dans un repère}
Dans un repère orthonormé, $A(0;0)$, $B(6;0)$ et $C(0;8)$, unité 1 cm. Calculer BC.
\label{fig-ex-13-9}\figureGeo{pyth-09}
\end{exercice}\end{blocLivre}
\begin{blocLivre}{ex-13-10}\begin{exercice}[2]{Trouver le bon sommet}
$AB=10{,}5$, $AC=6{,}3$, $BC=8{,}4$ cm. Démontrer que ABC est rectangle et préciser le sommet.
\label{fig-ex-13-10}\figureGeo{pyth-10}
\end{exercice}\end{blocLivre}

\rubrique{Problèmes}
\begin{blocLivre}{ex-13-11}\begin{exercice}[3]{Un losange}
Un losange a pour diagonales 10 et 24 cm. Elles se coupent perpendiculairement en leur milieu. Calculer le côté et le périmètre.
\label{fig-ex-13-11}\figureGeo{pyth-11}
\end{exercice}\end{blocLivre}
\begin{blocLivre}{ex-13-12}\begin{exercice}[3]{Retrouver une aire}
Un rectangle a une diagonale de 13 cm et une largeur de 5 cm. Calculer sa longueur puis son aire.
\label{fig-ex-13-12}\figureGeo{pyth-12}
\end{exercice}\end{blocLivre}
\begin{blocLivre}{ex-13-13}\begin{exercice}[3]{Deux diagonales}
Un câble suit la diagonale d’un rectangle de 8 m sur 6 m, puis celle d’un rectangle de 10 m sur 24 m. Les deux parcours se rejoignent sans longueur supplémentaire. Calculer la longueur totale.
\label{fig-ex-13-13}\figureGeo{pyth-13}
\end{exercice}\end{blocLivre}
\begin{blocLivre}{ex-13-14}\begin{exercice}[3]{Une précision décisive}
Une barre rigide de 8 cm doit relier les extrémités de deux côtés perpendiculaires de 4 cm et 7 cm. Est-elle assez longue ? Justifier exactement.
\label{fig-ex-13-14}\figureGeo{pyth-14}
\end{exercice}\end{blocLivre}
\begin{blocLivre}{ex-13-15}\begin{exercice}[3]{Choisir un trajet}
Un terrain rectangulaire mesure 12 m sur 5 m. Comparer un trajet le long de deux côtés adjacents avec la diagonale. Quelle longueur gagne-t-on avec la diagonale ?
\label{fig-ex-13-15}\figureGeo{pyth-15}
\end{exercice}\end{blocLivre}

\clearpage
